% Level ceiling (grade 5): speeds as "kilometres in one hour" in words,
% whole-number speeds, distance = speed x time with whole numbers only.
% No v = d/t with letters (literal expressions arrive in math grade 7),
% no \unit{km/h} compound notation (arrives grade 7 with quotients).
\chapter{速度：距离与时间}\label{ch:g5:speed-distance-time}

“我更快！”——“不，是我更快！”每个操场都用最实在的办法来分高下：
排好队，预备，跑！但怎样比较一位骑车的人，
和昨天在另一条路上跑步的人呢？
\omterm{def:g5:speed-distance-time:speed}{速度}正是物理用来比较那些永远无法并排同时进行的比赛的办法。

\section{比赛取胜的两种方式}

\begin{proposition}[公平比赛的规则]\label{prop:g5:speed-distance-time:rules}
比较快慢，有两种公平的办法：
\begin{enumerate}
  \item 跑\emph{相同的路程}时，\omterm{def:g5:speed-distance-time:speed}{速度}快的人需要
        \emph{更少的时间}——先冲过终点线的人赢得短跑；
  \item 在\emph{相同时间}内，\omterm{def:g5:speed-distance-time:speed}{速度}快的人跑过
        \emph{更长的路程}——跑满一分钟后，\omterm{def:g5:speed-distance-time:speed}{速度}快的游泳者
        只是游得更远。
\end{enumerate}
路程相同或时间相同：固定其中一个，比较另一个。
\emph{两者}都不同时比较——“我跑得更远！”“可你跑了一整个
下午！”——什么也说明不了。
\end{proposition}

\begin{example}[不公平与公平]\label{ex:g5:speed-distance-time:fair}
莉娜骑了 \qty{10}{km}；马可骑了 \qty{15}{km}。
马可是否更快？无从判断——也许他骑的时间是别人的三倍。
但如果两人都恰好骑了一\omterm{ex:g2:measuring-time:clock}{小时}，比较就变得公平了：
时间相同、路程更长，马可就更快。这个小小的“一\omterm{ex:g2:measuring-time:clock}{小时}内”
就要成为我们最喜欢的说法。
\end{example}

\section{速度：一小时内的路程}

\begin{definition}[速度]\label{def:g5:speed-distance-time:speed}
匀速前进者的\emph{速度}\index{速度}，是指它在一段固定时间内
所走过的距离——通常是一\omterm{ex:g2:measuring-time:clock}{小时}。
“这辆汽车每\omterm{ex:g2:measuring-time:clock}{小时}行驶 $60$ 千米”说的就是它的速度；
步行轻快的人每\omterm{ex:g2:measuring-time:clock}{小时}大约走 $5$ 千米，赛车骑手
每\omterm{ex:g2:measuring-time:clock}{小时}大约骑 $40$ 千米。对于爬得很慢的动物，我们换用更便于理解的\omterm{def:g2:measuring-length:measure}{单位}：一只
蜗牛每\omterm{ex:g2:measuring-time:clock}{小时}大约爬 $5$ 米。
\end{definition}

\begin{omfigure}
\begin{tikzpicture}[scale=0.85]
  \draw[-Stealth, very thick] (0,0) -- (11.2,0);
  \foreach \x/\name/\v in {0.9/蜗牛/{$5$ m}, 3.1/步行者/{$5$ km},
      5.3/骑车人/{$20$ km}, 7.5/汽车/{$90$ km},
      9.7/飞机/{$900$ km}} {
    \fill[omDef] (\x,0) circle (2.6pt);
    \node[above, font=\footnotesize] at (\x,0.25) {\name};
    \node[below, font=\footnotesize] at (\x,-0.25) {\v};
  }
  \node[below, font=\small] at (5.6,-0.9)
    {每位旅行者在一小时内所走的距离};
\end{tikzpicture}

{\small \omterm{def:g5:speed-distance-time:speed}{速度}阶梯：各位旅行者一\omterm{ex:g2:measuring-time:clock}{小时}内
所走的距离——从蜗牛的五米到客机的
九百千米。}
\end{omfigure}

\begin{method}[测量自己的速度]\label{met:g5:speed-distance-time:measure}
\begin{enumerate}
  \item \omterm{def:g2:measuring-length:measure}{测量}或找出一段恰好 \qty{100}{m} 长的路段（许多
        运动场的跑道上都有标记）；
  \item 以平常的步速走完，由朋友为你计时——
        假设用时一分钟；
  \item 一\omterm{ex:g2:measuring-time:clock}{小时}里有 $60$ 个这样的一分钟，所以在一\omterm{ex:g2:measuring-time:clock}{小时}内
        你将走过 $60 \times 100 = \qty{6000}{m}$；
  \item 用千米来表示：你的步行\omterm{def:g5:speed-distance-time:speed}{速度}是
        每\omterm{ex:g2:measuring-time:clock}{小时} $6$ 千米。
\end{enumerate}
改为跑步再测一次并比较——大多数孩子的\omterm{def:g5:speed-distance-time:speed}{速度}大约能翻一倍。
\end{method}

\section{预测路程}

\begin{proposition}[由速度求路程]\label{prop:g5:speed-distance-time:distance}
匀速前进者此后每一\omterm{ex:g2:measuring-time:clock}{小时}都会再走
一段与一\omterm{ex:g2:measuring-time:clock}{小时}路程相同的距离。于是有：
\[
\text{所走路程} = \text{一小时内的路程} \times
\text{小时数}.
\]
每\omterm{ex:g2:measuring-time:clock}{小时}行驶 $60$ 千米的汽车，在 $3$ \omterm{ex:g2:measuring-time:clock}{小时}内
走过的路程为 $60 \times 3 = \qty{180}{km}$。
\end{proposition}

\begin{example}[行程计算]\label{ex:g5:speed-distance-time:journeys}
步行者每\omterm{ex:g2:measuring-time:clock}{小时}走 $4$ 千米，走了 $2$ \omterm{ex:g2:measuring-time:clock}{小时}：
$4 \times 2 = \qty{8}{km}$。骑车者每\omterm{ex:g2:measuring-time:clock}{小时}骑 $15$ 千米，
骑了 $4$ \omterm{ex:g2:measuring-time:clock}{小时}：$15 \times 4 = \qty{60}{km}$。客机
每\omterm{ex:g2:measuring-time:clock}{小时}飞行 $900$ 千米，飞了 $5$ \omterm{ex:g2:measuring-time:clock}{小时}：$900 \times 5 =
\qty{4500}{km}$——两顿饭之间就能横穿一个大洲。
\end{example}

\begin{omfigure}
\begin{tikzpicture}
\begin{axis}[omaxis, width=10.5cm, height=5.6cm,
    xlabel={行驶时间（小时）}, ylabel={所走路程（千米）},
    xtick={0,1,2,3}, ytick={0,15,30,45,60},
    ymin=0, ymax=68, xmin=0, xmax=3.4]
  \addplot[very thick, omDef] coordinates
    {(0,0) (1,15) (2,30) (3,45)};
  \addplot[very thick, omThm] coordinates
    {(0,0) (1,4) (2,8) (3,12)};
  \node[right, font=\small, omDef] at (axis cs:2.05,33)
    {骑车者：每小时 $15$ 千米};
  \node[right, font=\small, omThm] at (axis cs:2.05,10.5)
    {步行者：每小时 $4$ 千米};
\end{axis}
\end{tikzpicture}

{\small 匀速行程用线段表示：每\omterm{ex:g2:measuring-time:clock}{小时}增加的路程都相同。
\omterm{def:g5:speed-distance-time:speed}{速度}较快者的线段上升得更
陡峭。}
\end{omfigure}

\begin{remark}[匀速是一种简化]\label{rem:g5:speed-distance-time:steady}
真实的行进很少是匀速的：汽车经过村庄时会减速，
骑车人下坡飞驰、上坡费力，你路过面包店时也会慢悠悠地走。
我们的算式假装整段行程都以同一个稳定的\omterm{def:g5:speed-distance-time:speed}{速度}进行——
这是一种有用的假设，对规划行程已经足够好。
至于如何诚实地描述由许多段\omterm{def:g5:speed-distance-time:speed}{速度}组成的行程，
以及“平均”一词真正的含义——
那是以后学年的内容，等你掌握了相应的数学再学。
\end{remark}

\begin{example}[认识身边的速度]\label{ex:g5:speed-distance-time:signs}
路牌和汽车仪表盘用的是“每\omterm{ex:g2:measuring-time:clock}{小时}多少千米”的说法——
圆形路牌上写着 $50$，就是告诉司机：每\omterm{ex:g2:measuring-time:clock}{小时}行驶最多 $50$ 千米，
这样的\omterm{def:g5:speed-distance-time:speed}{速度}能保护追着皮球跑的小孩。
高速列车的屏幕上可能显示 $300$；
山里的路标则更喜欢用\omterm{ex:g2:measuring-time:clock}{小时}而不是千米，
比如“湖：2 \omterm{ex:g2:measuring-time:clock}{小时}”——
相信行人自己知道自己的\omterm{def:g5:speed-distance-time:speed}{速度}。
\end{example}

\section{练习}

\begin{exercise}[$\star$]\label{exo:g5:speed-distance-time:1}
请说出比较快慢的两种公平方法。为什么只凭“我比你跑得远”，
就什么也证明不了？
\end{exercise}

\begin{exercise}[$\star$]\label{exo:g5:speed-distance-time:2}
两名游泳者游过同一座泳池：安娜需要 $30$ 秒，
博需要 $25$ 秒。谁更快？依据哪条公平比赛规则？
\end{exercise}

\begin{exercise}[$\star$]\label{exo:g5:speed-distance-time:3}
行驶一\omterm{ex:g2:measuring-time:clock}{小时}后，拖拉机走了 \qty{20}{km}，
踏板车走了 \qty{45}{km}。谁更快？依据哪条规则？
\end{exercise}

\begin{exercise}[$\star$]\label{exo:g5:speed-distance-time:4}
“这列火车每\omterm{ex:g2:measuring-time:clock}{小时}行驶 $200$ 千米”是什么意思？
它行驶 $2$ \omterm{ex:g2:measuring-time:clock}{小时}能走多远？半\omterm{ex:g2:measuring-time:clock}{小时}呢？
\end{exercise}

\begin{exercise}[$\star$]\label{exo:g5:speed-distance-time:5}
按\omterm{def:g5:speed-distance-time:speed}{速度}排序：步行者（每\omterm{ex:g2:measuring-time:clock}{小时} $5$ 千米）；\ 鸽子
（每\omterm{ex:g2:measuring-time:clock}{小时} $80$ 千米）；\ 蜗牛（每\omterm{ex:g2:measuring-time:clock}{小时} $5$ \emph{米}
）；\ 城市公交车（每\omterm{ex:g2:measuring-time:clock}{小时} $30$ 千米）。
\end{exercise}

\begin{exercise}[$\star$]\label{exo:g5:speed-distance-time:6}
你沿着做好标记的 \qty{100}{m} 路段散步，用了 $2$ 分钟。
照这个\omterm{def:g5:speed-distance-time:speed}{速度}，一\omterm{ex:g2:measuring-time:clock}{小时}能走多少米？
请用每\omterm{ex:g2:measuring-time:clock}{小时}多少千米来表示你的\omterm{def:g5:speed-distance-time:speed}{速度}。
\end{exercise}

\begin{exercise}[$\star$]\label{exo:g5:speed-distance-time:7}
一辆汽车保持每\omterm{ex:g2:measuring-time:clock}{小时} $80$ 千米的\omterm{def:g5:speed-distance-time:speed}{速度}，持续行驶 $3$ \omterm{ex:g2:measuring-time:clock}{小时}。
它一共走了多远？再看骑车人以每\omterm{ex:g2:measuring-time:clock}{小时} $15$ 千米的\omterm{def:g5:speed-distance-time:speed}{速度}，
在同样的 $3$ \omterm{ex:g2:measuring-time:clock}{小时}内又走了多远？
\end{exercise}

\begin{exercise}[$\star$]\label{exo:g5:speed-distance-time:8}
在本章的行程图上，怎样一眼看出\omterm{def:g5:speed-distance-time:speed}{速度}较快的人？
请从图上读出：骑车人
行驶 $2$ \omterm{ex:g2:measuring-time:clock}{小时}后走了多远？
\end{exercise}

\begin{exercise}[$\star\star$]\label{exo:g5:speed-distance-time:9}
一艘渡船每\omterm{ex:g2:measuring-time:clock}{小时}航行 $30$ 千米，渡河全程为
\qty{90}{km}。利用
\cref{prop:g5:speed-distance-time:distance} 反向使用——试着
逐个试算\omterm{ex:g2:measuring-time:clock}{小时}数，渡河需要多长时间？
\end{exercise}

\begin{exercise}[$\star\star$]\label{exo:g5:speed-distance-time:10}
两座城镇相距 \qty{60}{km}。玛拉从其中一座城镇出发，\omterm{def:g5:speed-distance-time:speed}{速度}为每\omterm{ex:g2:measuring-time:clock}{小时} $20$
千米；同一时刻，乔恩从另一座城镇出发，
朝她骑来，\omterm{def:g5:speed-distance-time:speed}{速度}为每\omterm{ex:g2:measuring-time:clock}{小时} $10$ 千米。两人之间的距离
每\omterm{ex:g2:measuring-time:clock}{小时}缩短多少千米？
经过多少\omterm{ex:g2:measuring-time:clock}{小时}后两人相遇？
\end{exercise}

\begin{exercise}[$\star\star$]\label{exo:g5:speed-distance-time:11}
这是一则古老的寓言，现在配上数字：赛道全长 \qty{120}{m}。
兔子每分钟跑 $60$ 米——但跑了一分钟之后
就停下来睡一个 $70$ 分钟的觉，然后继续跑。乌龟
每分钟慢慢爬 $2$ 米，从不停下。兔子开始午睡时
在哪里？两者各在什么时刻到达终点线？
谁会赢？
\end{exercise}
